\documentclass[./测控三级项目.tex]{subfiles} % 使用 subfiles 文档类，指定主文档

\begin{document}
	\subsection{题目重述}
	\begin{compactenum}
		\item 查阅资料（参考课本§1.8），简述运算放大器电源供电电路结构原理。
		\item 选用合适的接口器件和电路元件正确焊装，在引出排针上可测到输入的电源电压（万用表）。
		\item 会使用线性稳压电源，能同时给电路提供要求的正负电源。\label{aa}
	\end{compactenum}
	\subsection{题目分析与实施方案}
	\subsubsection{第一问}
	本题要求通过查阅课本 §1.8 资料，阐述运算放大器电源供电电路的结构原理，包括电源类型、电路元件作用及其连接方式等理论知识。\par
	从课本 §1.8 获取运算放大器电源供电电路相关内容，包括电路图、文字讲解等，了解双电源、单电源供电方式，去耦电容等元件的作用。此步骤满足题目对查阅资料和简述原理的要求。\par
	整理资料内容，明确电路结构组成部分，如正电源（$V_{CC}$）、负电源（$V_{EE}$）的作用，不同电源供电下运算放大器的工作情况，以及去耦电容在减少电源噪声方面的原理。这一步深化对原理的理解和阐述。
	\subsubsection{第二问}
	根据作业指导书所给的原理图，选用$5.08$mm螺钉式接线端子作为电源输入端，同时在正电源和地、负电源和地之间并联$10$uF和$100$nFMLCC电容作为电源滤波。\par
	实际使用公对公杜邦线和鳄鱼夹将线性直流电源输出口与PCB电源网络相连。
	\subsubsection{第三问}
	将线性稳压电源的交流输入插头连接到市电插座。根据运算放大器供电要求，调节稳压电源的正电压输出通道旋钮，设置为所需正电源电压（如+15V），并连接到运算放大器的正电源引脚；同样，调节负电压输出通道旋钮设置负电源电压（如-15V），连接到负电源引脚。
	\subsection{理论分析与仿真}
	\subsubsection{第一问}
	\textbf{双电源供电原理}\quad 运算放大器使用双电源供电是为了能够有效处理既有正电压又有负电压的信号。双电源由正电源（$V_{CC}$）和负电源（$V_{EE}$）组成，以地（GND）为参考电位。当输入信号为正时，正电源提供能量使运算放大器能够正常放大信号并输出正电压；当输入信号为负时，负电源提供能量使运算放大器输出负电压。这样可以确保运算放大器能够对输入信号的整个动态范围进行放大，避免信号失真。\par
	\textbf{电源幅值}\quad 运算放大器的电源幅值大小影响了输出信号的最大摆幅。通常$\pm15$V数值长期作为模拟系统的标准，但是也有如$\pm12$V或$\pm5$V等双电源供电方案。在简单的模拟电路中，通常使用单电源输入，去除负电压部分，因此电源负极直接和地相接。\par
	\textbf{单电源供电对电路性能的影响}\quad 单电源供电时，由于没有负电源，运算放大器无法输出负电压，这使得输入信号的负半周无法得到正常放大，导致信号动态范围减小。\par
	\textbf{单电源供电的特殊器件与电路}\quad 为了使运算放大器在单电源供电下能够正常处理信号，需要设置偏置电路。通常是通过电阻分压网络来为运算放大器的同相输入端提供一个合适的直流偏置电压，使输入信号能够在这个偏置电压的基础上进行放大。\par
	\textbf{电解电容和瓷片电容差别}\quad 电解电容和瓷片电容在材料结构、电容量范围、频率特性、耐压特性等问题上有所区别，详细区别详见表\ref{tab:电解电容和瓷片电容差别}。
	\begin{table}[h!]
		\centering
		\caption{电解电容和瓷片电容差别}
		\label{tab:电解电容和瓷片电容差别}
		\begin{tabular}{|p{2cm}|p{7cm}|p{7cm}|}
			\hline
			比较项目 & 电解电容  & 瓷片电容 \\
			\hline
			材料与结构 & 由阳极金属箔（通常为铝）、阴极电解质和隔离纸组成，有极性。阳极金属箔表面有氧化膜作为电介质，电解液参与电荷存储过程，整体结构相对松散，体积较大，正负极标识明显。 & 以陶瓷材料为介质，电极一般为银等金属材料。多层瓷片电容是多层陶瓷介质与金属电极交替叠加而成，无极性，结构紧凑，体积小，外观上无明显正负极标识（特殊情况除外）。 \\
			\hline
			电容量范围 & 电容量范围大，从几微法到数千微法甚至更大。通过增大电极面积和调整电介质厚度等方法可实现较大电容值，能满足低频大电容需求电路。 & 电容量通常较小，一般在几皮法到零点几微法，虽有大容量瓷片电容但不常见。电容值主要取决于陶瓷介质介电常数、电极面积和介质厚度，高频应用决定其电容值不会太大。 \\
			\hline
			频率特性 & 高频性能差，等效串联电感（ESL）和等效串联电阻（ESR）较大。电解液和电极结构导致其对高频信号响应慢，在高频电路中会因 ESL 产生较大感抗、ESR 产生电压降和功率损耗，滤波效果差且可能引入高频噪声。 & 高频特性好，ESL 和 ESR 小。陶瓷介质和紧凑结构使它能快速响应高频信号变化，在高频下有效滤除高频信号，保障高频电路信号传输质量。 \\
			\hline
			耐压特性 & 耐压值有一定范围，可较高，从十几伏到几百伏都有。耐压能力与电介质厚度、质量及电解液成分等有关，能用于多种电压等级电路，尤其在高电压低频电源电路中常用。& 耐压值一般较低，多在几十伏以内，少数特殊瓷片电容耐压值可达较高水平。耐压取决于陶瓷介质击穿强度和电极间距等，常用于低电压高频电路。 \\
			\hline
			稳定性与寿命&	稳定性较差，电解液会随时间、温度变化，在高温、高湿度等恶劣环境下，电解液可能干涸、变质，导致电容参数改变，如电容量下降、ESR 增大，影响电路性能，寿命受环境因素影响大。&	稳定性好，陶瓷介质和金属电极化学性质稳定，受温度、时间影响小，在正常使用条件下能长期保持性能稳定，参数变化小，使用寿命长。\\
			\hline
			价格与应用场景&	价格相对较低，特别是普通规格的铝电解电容。由于成本低，常用于对电容精度要求不高、注重成本控制的低频滤波、电源去耦电路，如普通电子设备电源输入输出端的滤波。&	价格因容量、精度、材质等而异，小容量普通精度的价格较低，高精度或特殊材质的价格较高。广泛用于高频滤波、信号耦合等高频电路，如射频电路中的信号处理、数字电路中高频时钟信号的噪声滤除。 \\
			\hline
		\end{tabular}
	\end{table}	\par
	\textbf{铝电解电容和钽电容差别}\quad 铝电解电容和钽电容在材料结构、电容量范围、频率特性、耐压特性等问题上有所区别，详细区别详见表\ref{tab:铝电解电容和钽电容差别}。
	\begin{table}[H]
		\centering
		\caption{铝电解电容和钽电容差别}
		\label{tab:铝电解电容和钽电容差别}
		\begin{tabular}{|p{2.5cm}|p{7cm}|p{7cm}|}
			\hline
			比较项目 & 电解电容  & 瓷片电容 \\
			\hline
			材料与结构 & 以铝箔作为阳极，铝箔表面通过电化学方法形成一层氧化铝作为电介质。阴极是浸有电解液的纸或其他多孔材料，电解液在电场作用下发生离子移动来存储电荷。整体结构相对简单，成本较低。铝电解电容通常为圆柱形，体积较大，并且有明显的正负极标识。 & 以钽金属作为阳极，通过在钽表面形成一层钽的氧化物作为电介质。电解质可以是固体（如二氧化锰）也可以是液体，采用固体电解质的钽电容密封性更好。钽电容的结构更加紧凑，形状多样，常见的有贴片式和圆柱形。其内部结构设计使得它在相同电容量下体积更小，也有正负极之分，不过部分钽电容的正负极标识可能不太明显。\\
			\hline
			电容量范围 & 电容量范围非常广泛，从几微法到数千微法都有。这是因为铝箔电极可以有较大的面积，并且可以通过调整电解液的成分和电介质的厚度等因素来改变电容量。能够满足很多需要大电容值的电路。 & 电容量范围一般在零点几微法到几百微法之间，容量相对较小。这主要是由钽电容的材料和结构特点决定的，尽管如此，在一些对电容体积要求较小但对性能要求较高的电路中，其容量也能满足需求。\\
			\hline
			频率特性 & 高频性能较差。由于其内部存在电解液，导致等效串联电感（ESL）和等效串联电阻（ESR）较大。在高频环境下，ESL 会产生较大的感抗，阻碍电流的快速变化，ESR 会导致明显的电压降和功率损耗，所以对高频信号的滤波效果不好。在一些高频电路中，铝电解电容可能会引入额外的高频噪声，影响电路性能。 & 具有较好的高频特性。其 ESL 和 ESR 相对较小，这是因为钽电容的结构和材料使得它在高频下能够更快地响应电流变化。在高频电路中，钽电容能够更有效地滤除高频信号，减少信号失真，例如在射频电路的滤波和信号耦合环节表现更好。 \\
			\hline
			耐压特性 & 耐压值有多种选择，一般可以达到较高的值，通常从几伏到几百伏不等。耐压能力主要取决于电介质的厚度和质量，以及电解液的成分等因素。在高电压的电源电路中应用较为广泛，如工业电源、电力电子设备中的滤波电路等。& 耐压值相对较低，通常在几伏到几十伏之间。这是因为钽电容的材料和结构特点限制了其承受高电压的能力。在高电压应用场景下，钽电容可能会发生击穿损坏，所以主要用于低电压、高性能要求的电路。 \\
			\hline
			稳定性与寿命&	稳定性稍差。其电解液会随着时间和温度的变化而逐渐干涸或变质，尤其是在高温环境下，这种变化更为明显。当电解液性能下降时，电容的电容量、ESR 等参数都会发生变化，导致电容性能下降，寿命相对较短。一般来说，其寿命受使用环境温度、工作电压、工作频率等因素影响较大。 &	稳定性好。钽电容受温度和时间的影响较小，其电容量和 ESR 等参数在长时间使用和不同温度环境下能保持相对稳定。这是因为钽的氧化物介质和固体电解质（如果是固体钽电容）具有较好的化学稳定性，使得钽电容具有较长的使用寿命，能够在较为恶劣的环境下保持良好的性能。\\
			\hline
			价格与应用场景&价格便宜。由于原材料成本较低，生产工艺相对简单，所以铝电解电容价格较低。常用于对电容精度要求不高、成本敏感的低频滤波、电源去耦等电路。例如在普通的电子产品电源电路中，用于平滑直流电压的低频纹波，或者在音频放大电路中用于耦合和旁路低频信号。&	价格较高。钽是一种稀有金属，且钽电容的生产工艺要求较高，这使得其成本上升。主要用于对电容性能要求较高的高精度、高频滤波、小型化的电子设备中。例如在一些高端的通信设备、精密仪器、航空航天电子设备等中，用于滤除高频噪声、保证信号质量和电路的稳定性。 \\
			\hline
		\end{tabular}
	\end{table}			
	\subsection{实际电路测试结果与数据分析}
	\subsubsection{第二问}
	根据《测控仪器电路项目作业指导书》附件五给出的电路板原理图，可以得知，在P0处应当焊接一个螺钉式连接器，后续焊接四个容值为100nF的MLCC电容，根据焊盘尺寸，两个为1206封装，两个为0805封装。并焊接两个4P排针，作为正负电源的引出。\par
	由于在实验室中未找到封装为1206的100nF电容，因此四个电容全部使用0805封装。
	其焊接完成图如图\ref{fig:电源部分焊接完成}所示。
	\begin{figure}[h]
		\centering
		\includegraphics[width=0.45\textwidth]{电源部分焊接完成}
		\caption{电源部分焊接完成照片}
		\label{fig:电源部分焊接完成}
	\end{figure}\par
	焊接完成后，使用万用表蜂鸣器档测试电路通断，测试发现排针与螺钉式连接器的螺钉短路，蜂鸣器鸣叫，因此证明电路电源输入部分通断无误，电压测量在第三问处上电后才能测量，因此具体测量现象将在第三问给出。
	\subsubsection{第三问}
	本文要求上电测试，因此需要使用线性直流电源进行供电。首先打开线性直流电源的总电源，将电压通道CH1和CH2电压值都调整到15V。为了防止误触等原因在上电之后发生短路，将两个电流限制在0.1A左右，以防止芯片的烧毁。\par
	参数修改完成后，线性电源面板如图\ref{fig:直流电源未开启状态}所示。
	\begin{figure}[H]
		\centering
		\includegraphics[width=0.4\textwidth]{直流电源未开启状态}
		\caption{直流电源未开启状态时面板照片}
		\label{fig:直流电源未开启状态}
	\end{figure}\par
	实验统一将通道CH1认为是正电源（15V）供电，将通道CH2认为是负电源（-15V）供电，因此将CH1的电源负输出端和CH2的正输出端短接，接入电路板的GND；CH1的电源正输出端接入电路板的$V_{CC}$；CH2的电源负输出端接入电路板的$V_{EE}$。\par
	电源接入后，点击电源输出按键。电路输出后，电源面板如图\ref{fig:直流电源开启状态}所示。
	\begin{figure}[H]
		\centering
		\includegraphics[width=0.4\textwidth]{直流电源开启状态}
		\caption{直流电源开启状态时面板照片}
		\label{fig:直流电源开启状态}
	\end{figure}\par
	面板显示电流分别为0.05A和0.06A，电压维持15V没有变化，说明没有限流，电路正常。\par
	通过使用电压表测量GND与$V_{CC}$和$V_{EE}$之间的电压，如图\ref{fig:电源电压测量照片}所示。正电源电压约为14.98V，负电源电压约为-14.95V。这与上述所设置的电压差距不大，可认为电源部分调试正确。
	\begin{figure}[h]
		\begin{subfigure}{0.5\linewidth}
			\centering
			\includegraphics[width=0.8\linewidth]{测量正电源电压}
			\caption{正电源电压测量照片}
			\label{fig:测量正电源电压}
		\end{subfigure}
		\begin{subfigure}{0.5\linewidth}
			\centering
			\includegraphics[width=0.8\linewidth]{测量负电源电压}
			\caption{负电源电压测量照片}
			\label{fig:测量负电源电压}
		\end{subfigure}
		\caption{电源电压测量照片}
		\label{fig:电源电压测量照片}
	\end{figure}
	
	\subsection{结论与收获}
	\textbf{结论}：通过第一题我们了解了运算放大器电源供电电路的结构原理，包括电源类型、电路元件作用及其连接方式等理论知识，还调试出来了稳定的供电电路，为之后的问题开了一个好头。\par
	\textbf{收获}：这是我们小组的第一次合作，大家都十分积极的讨论，工作开展的也特别愉快。郭嘉媛对供电电路有了更加深刻的了解，自己动手焊接虽然不是十分完美，但也较为出色的完成了任务。在线路调试时，游寅龙在小差错中则明白了不能一味相信经验，要大胆尝试新方法，而郑嘉宁则明白了不会没关系，只要肯学就一定会成功。
	
\end{document}